\section*{Capítulo \ref{ch:b2:atmospheric-thermodynamics} --- Termodinâmica da atmosfera}

\begin{solution}{exo:b2:atmospheric-thermodynamics:1}
(a) $e = e_s(\qty{278.15}{K}) = \qty{8.73}{hPa}$ e $e_s(\qty{283.15}{K}) = \qty{12.32}{hPa}$.
(b) $8.73/12.32 = \qty{70.9}{\percent}$. (c) $r = 0.622 \times 8.73/(850 - 8.73) =
\qty{6.45}{g/kg}$ e $q = r/(1 + r) = \qty{6.41}{g/kg}$.
\end{solution}

\begin{solution}{exo:b2:atmospheric-thermodynamics:2}
(a) $\mathrm{d}\ln e_s/\mathrm{d}T = 5419/T^2$: \qty{7.26}{\percent} por kelvin a
\qty{273.15}{K}, \qty{6.31}{\percent} a \qty{293.15}{K}, \qty{5.90}{\percent} a
\qty{303.15}{K}. (b) $e_s(\qty{20}{\celsius}) = \qty{23.66}{hPa}$ e $e_s(\qty{25}{\celsius})
= \qty{32.26}{hPa}$; razão 1.363, contra $1.063^5 = 1.357$. A taxa é uma derivada de
$\ln e_s$, de modo que em \qty{5}{K} o fator é $\exp(5 \times 0.0620) = 1.363$, sendo 0.0620 a
taxa no ponto médio, \qty{22.5}{\celsius}; compor \qty{6.3}{\percent} cinco vezes chega perto,
mas não é exato. (c) $23.66/23.39 = 1.012$: a fórmula dá \qty{1.2}{\percent} a mais, porque
conserva o calor latente do ponto triplo, enquanto o calor latente diminui quando a temperatura
sobe; a inclinação da fórmula é forte demais, e seu erro cresce com a distância a
\qty{0}{\celsius}.
\end{solution}

\begin{solution}{exo:b2:atmospheric-thermodynamics:3}
(a) $125 \times (30 - 14) = \qty{2000}{m}$. (b) $30 - 9.8 \times 2.0 = \qty{10.4}{\celsius}$.
(c) $14 - 1.8 \times 2.0 = \qty{10.4}{\celsius}$: a temperatura e o ponto de orvalho se encontram, como devem no
nível de saturação; a regra é esse cálculo. (d) $125 \times (33 - 14) = \qty{2375}{m}$:
as bases sobem à medida que a tarde aquece o ar junto ao solo.
\end{solution}

\begin{solution}{exo:b2:atmospheric-thermodynamics:4}
(a) A curva do ponto de orvalho coincide com a curva da temperatura do solo até cerca de \qty{908}{hPa}:
o meio quilômetro mais baixo está saturado, ou quase. (b) Entre cerca de 908 e \qty{889}{hPa}
a curva da temperatura volta-se para a direita, para valores mais quentes, à medida que a pressão diminui: a
temperatura sobe com a altitude. (c) NCL a cerca de \qty{964}{hPa}, algumas dezenas de metros acima; NCE a
cerca de \qty{730}{hPa}; NE a cerca de \qty{235}{hPa}. (d) No topo da inversão o ponto de orvalho
está \qty{11}{K} abaixo da temperatura, e \qty{20}{K} abaixo dela a \qty{850}{hPa}: um ar quente e muito seco
assenta-se sobre um ar fresco e úmido. Sua umidade relativa é baixa, de modo que ele não adquiriu seu
calor por condensação; um ar que desce aquece-se por compressão segundo o gradiente adiabático seco
conservando sua \omterm{def:b2:atmospheric-thermodynamics:moist-air}{razão de mistura}, o que daria uma camada assim.
\end{solution}

\begin{solution}{exo:b2:atmospheric-thermodynamics:5}
(a) $T_v = 303.15 \times (1 + 0.020/0.622)/1.020 = \qty{306.76}{K}$; a aproximação dá
$303.15 \times (1 + 0.61 \times 0.020) = \qty{306.85}{K}$. (b) $\rho = p/(R_dT_v) = 10^5/(287.0
\times 306.76) = \qty{1.1358}{kg/m^3}$, contra $10^5/(287.0 \times 303.15) =
\qty{1.1494}{kg/m^3}$ para o ar seco: \qty{1.2}{\percent} mais leve. (c) Em $T_v - T =
\qty{3.6}{K}$. (d) Uma parcela que sobe só precisa ser mais leve que o ambiente: sobre mares quentes
a umidade do ar acrescenta um empuxo equivalente a vários kelvin, tanto quanto um excesso
apreciável de temperatura.
\end{solution}

\begin{solution}{exo:b2:atmospheric-thermodynamics:6}
(a) $e_s$ = 6.11, 23.66 e \qty{43.54}{hPa}; $r_s = 0.622\,e_s/(1000 - e_s)$ = 3.82, 15.07 e
\qty{28.31}{g/kg}; $\Gamma_s$ = 6.48, 4.20 e \qty{3.44}{K/km}. (b) A parcela seca resfria-se
\qty{9.8}{K}, a saturada \qty{4.2}{K}: a parcela saturada está \qty{5.6}{K} mais quente.
(c) Mais acima o ar é mais frio, $r_s$ é menor, resta pouco vapor para condensar, e
$\Gamma_s$ volta a se aproximar de $\Gamma_d$.
\end{solution}

\begin{solution}{exo:b2:atmospheric-thermodynamics:7}
(a) $e = e_s(\qty{291.65}{K}) = \qty{21.52}{hPa}$, $r = 0.622 \times 21.52/(971.1 - 21.52) =
\qty{14.1}{g/kg}$; $\theta = 292.15 \times (1000/971.1)^{0.2857} = \qty{294.6}{K}$; $\theta_e =
294.6 \exp[2.501 \times 10^6 \times 0.0141/(1004.5 \times 291.5)] = 294.6 \times 1.128 =
\qty{332.3}{K}$. (b) $e_s(\qty{264.25}{K}) = \qty{3.13}{hPa}$, $r_s = 0.622 \times 3.13/(500 -
3.13) = \qty{3.92}{g/kg}$; $\theta = 264.25 \times 2^{0.2857} = \qty{322.1}{K}$; a parcela está
saturada, logo $T_{\text{NCL}} = T$ e $\theta_e = 322.1 \exp[2.501 \times 10^6 \times
0.00392/(1004.5 \times 264.25)] = \qty{334.3}{K}$. (c) $\theta$ subiu \qty{27.5}{K}, aquecida
pelo calor latente dos \qty{10.2}{g/kg} de vapor que condensaram; $\theta_e$ variou
\qty{0.6}{\percent}, dentro da precisão de sua fórmula aproximada: ela se conserva.
\end{solution}

\begin{solution}{exo:b2:atmospheric-thermodynamics:8}
(a) $\text{CAPE} = 287.0 \times 4 \times \ln(700/300) = 287.0 \times 4 \times 0.847 =
\qty{973}{J/kg}$. (b) $\sqrt{2 \times 973} = \qty{44.1}{m/s}$. (c) Entre 1 e
\qty{1500}{J/kg}: positiva, mas não grande. (d) Uma faixa de largura horizontal constante entre duas
curvas; pela \cref{prop:b2:atmospheric-thermodynamics:areas} sua área é proporcional a $\int
\Delta T\,\mathrm{d}\ln p = 4 \ln(700/300)$, qualquer que seja a inclinação das isotermas.
\end{solution}

\begin{solution}{exo:b2:atmospheric-thermodynamics:9}
(a) $D_c = \sqrt{3 \times 0.61 \times 80^3/2.1}\,\unit{nm} = \qty{668}{nm} =
\qty{0.668}{\micro m}$; $s_c = 2A/(3D_c) = 2 \times 2.1/(3 \times 668) = 0.00210 =
\qty{0.210}{\percent}$. (b) Resolvendo $s_c = \sqrt{4A^3/(27\kappa D_d^3)}$ em $D_d$: $D_d =
\bigl(4A^3/(27\kappa s_c^2)\bigr)^{1/3}$, \qty{63}{nm} a \qty{0.3}{\percent} e \qty{131}{nm} a
\qty{0.1}{\percent}. (c) Todas as partículas entre 63 e \qty{131}{nm} passam a se ativar também: mais
gotículas e, para a mesma água líquida, menores.
\end{solution}

\begin{solution}{exo:b2:atmospheric-thermodynamics:10}
(a) $w = \sqrt{2 \times 130} = \qty{16.1}{m/s}$, impossível para um ar em repouso junto ao solo.
(b) A temperatura potencial do topo da inversão é $293.65 \times (1000/889.1)^{0.2857} =
\qty{303.7}{K}$; uma parcela da superfície com esse $\theta$ tem $T = 303.7 \times (971.1/1000)^{0.2857}
= \qty{301.2}{K}$, \qty{28.0}{\celsius}. (c) $125 \times (28.0 - 18.5) = \qty{1190}{m}$ acima do
solo, acima do antigo topo da inversão. (d) Uma frente, que levanta o ar sobre sua rampa; a lenta ascensão
do ar à frente de uma perturbação em altitude, que resfria a própria tampa.
\end{solution}

\begin{solution}{exo:b2:atmospheric-thermodynamics:11}
(a) Quando $r_s \to 0$, o numerador tende a $g$ e o denominador a $c_p$. (b) Escreva-se
$L^2r_s/(R_vT^2) = c_p \times \bigl(Lr_s/(R_dT)\bigr) \times LR_d/(R_vTc_p)$, com
$LR_d/(R_vTc_p) = 5.1$ a \qty{303}{K}: quando $Lr_s/(R_dT) \gg 1$, também $L^2r_s/(R_vT^2) \gg c_p$,
e $\Gamma_s \to g\,(Lr_s/(R_dT))/(L^2r_s/(R_vT^2)) = gR_vT/(LR_d) = 9.81 \times 461.5 \times
303/(2.501 \times 10^6 \times 287.0) = \qty{1.91}{K/km}$. (c) Com $T_d = T$,
$\Gamma_{\text{orv}} = gR_vT/(LR_d)$, o mesmo. Nesse limite o calor latente é tão grande que
a parcela quase não pode condensar seu vapor sem se aquecer muito acima de seu ponto de
saturação: ela conserva sua \omterm{def:b2:atmospheric-thermodynamics:moist-air}{razão de mistura} e segue a linha de \omterm{def:b2:atmospheric-thermodynamics:moist-air}{razão de mistura} constante, ao longo da qual a
temperatura cai segundo $\Gamma_{\text{orv}}$. (d) A \qty{30}{\celsius} e \qty{1000}{hPa},
$Lr_s/(R_dT) = 0.81$, não muito maior que 1, e $L^2r_s/(R_vT^2) = 4180$, quatro vezes $c_p$:
$\Gamma_s = \qty{3.4}{K/km}$, ainda longe de \qty{1.9}{K/km}. O limite não é atingido.
\end{solution}

\begin{solution}{exo:b2:atmospheric-thermodynamics:12}
(a) $2.868/2.599 = 1.104$ e $1.257/1.032 = 1.218$: \qty{10.4}{\percent} e
\qty{21.8}{\percent}. (b) As gotículas ativam-se a alguns décimos de por cento de \omterm{def:b2:atmospheric-thermodynamics:microphysics}{supersaturação},
dezenas de vezes menos. (c) Os cristais crescem depressa por deposição; retiram vapor do ar,
que fica abaixo da saturação em relação às gotículas; as gotículas evaporam e seu vapor vai para
o gelo (o processo de Bergeron), até a nuvem estar glaciada. (d) A \qty{0}{\celsius} as duas
\omterm{def:b2:atmospheric-thermodynamics:moist-air}{pressões de vapor} de saturação são iguais: encontram-se no ponto triplo, e a razão tende a um
perto dele.
\end{solution}

\begin{solution}{pb:b2:atmospheric-thermodynamics:1}
\textbf{1.} $e = e_s(\qty{291.65}{K}) = \qty{21.5}{hPa}$.
\textbf{2.} $r = 0.622 \times 21.5/(971.1 - 21.5) = \qty{14.1}{g/kg}$.
\textbf{3.} $q = 14.1/1.0141 = \qty{13.9}{g/kg}$.
\textbf{4.} $e_s(\qty{292.15}{K}) = \qty{22.2}{hPa}$; $21.5/22.2 = \qty{97}{\percent}$.
\textbf{5.} $0.61 \times 0.0141 \times 292.15 = \qty{2.5}{K}$.
\textbf{6.} $125 \times 0.5 = \qty{63}{m}$; $p = 971.1 \exp[-9.81 \times 62.5/(287.0 \times 294.7)]
= \qty{964}{hPa}$.
\textbf{7.} $\theta = 292.15 \times (1000/971.1)^{0.2857} = \qty{294.6}{K}$.
\textbf{8.} $16.9 - 17.3 = \qty{-0.4}{K}$: mais fria.
\textbf{9.} $(20.5 - 16.4)/(1102 - 924) = 4.1/178$: a temperatura sobe \qty{23}{K/km}, um
gradiente de \qty{-23}{K/km}. Uma inversão: absolutamente estável, uma tampa.
\textbf{10.} $15.4 - 20.5 = \qty{-5.1}{K}$.
\textbf{11.} $e_s(\qty{-1.8}{\celsius}) = \qty{5.4}{hPa}$ e $e_s(\qty{18.6}{\celsius}) =
\qty{21.7}{hPa}$: \qty{25}{\percent}. Quente e seco, pode ser um ar que desceu e se aqueceu por
compressão, ou um ar aquecido sobre um solo seco a barlavento; não é o ar úmido de baixo.
\textbf{12.} $18.6 + 18.4 = \qty{37.0}{K}$ em $5725 - 1489 = \qty{4236}{m}$, \qty{8.73}{K/km}: entre $\Gamma_s$ e
$\Gamma_d$, condicionalmente instável.
\textbf{13.} $-18.4 - (-8.9) = \qty{-9.5}{K}$: instabilidade extrema (de $-8$ a $-10$).
\textbf{14.} Do solo para cima: 0, $-0.4$, $-0.2$, $-5.1$, $-4.9$, $+1.6$, $+9.5$, $+10.0$,
$+5.6$, $+1.5$, $-9.5$~\unit{K}.
\textbf{15.} Entre 850 ($-4.9$) e \qty{700}{hPa} ($+1.6$), na fração $4.9/6.5 = 0.754$:
$\ln p = \ln 850 + 0.754 \ln(700/850)$, $p = \qty{734}{hPa}$; $z = 1489 + 0.754 \times 1621 =
\qty{2711}{m}$, \qty{2.37}{km} acima do solo.
\textbf{16.} Entre 250 ($+1.5$) e \qty{200}{hPa} ($-9.5$), na fração $1.5/11.0 =
0.136$: $p = \qty{242.5}{hPa}$ e $z = 10560 + 0.136 \times 1446 = \qty{10757}{m}$, cerca de
\qty{200}{m} abaixo da tropopausa indicada (\qty{10954}{m}).
\textbf{17.} Os segmentos ($\Delta T$ médio vezes o $\ln$ da razão de pressões): $0.80 \times 0.0474
= 0.038$; $5.55 \times 0.3365 = 1.867$; $9.75 \times 0.2231 = 2.176$; $7.80 \times 0.2877 =
2.244$; $3.55 \times 0.1823 = 0.647$; $0.75 \times 0.0305 = 0.023$; soma \qty{6.995}{K}.
$\text{CAPE} = 287.0 \times 6.995 = \qty{2010}{J/kg}$, com três algarismos.
\textbf{18.} Entre 1500 e \qty{2500}{J/kg}: grande.
\textbf{19.} $-0.2 \times 0.0486 - 0.3 \times 0.0189 - 2.65 \times 0.0207 - 5.0 \times 0.0450 -
2.45 \times 0.1467 = -0.655$~\unit{K}; $\text{CIN} = 287.0 \times (-0.655) = \qty{-188}{J/kg}$.
\textbf{20.} A tabela guarda apenas onze níveis: entre eles a regra dos trapézios substitui as
curvas por segmentos retos em $\ln p$ e perde suas curvaturas. A grade de \qty{1}{hPa} segue as curvas. A \omterm{def:b2:atmospheric-thermodynamics:virtual}{temperatura virtual}
acrescenta o empuxo do vapor de água da parcela, mais abundante na parcela saturada que no
ar seco à sua volta: uma CAPE maior, uma CIN menor.
\textbf{21.} $\sqrt{2 \times 2010} = \qty{63}{m/s}$.
\textbf{22.} $10757 - 2711 = \qty{8046}{m}$ a $63.4/2 = \qty{31.7}{m/s}$: \qty{254}{s}, cerca de quatro
minutos.
\textbf{23.} O peso da água condensada que a corrente ascendente carrega; a mistura com o ar mais seco
à sua volta (entranhamento), que a resfria quando as gotículas evaporam.
\textbf{24.} A teoria da parcela dá à corrente ascendente de uma tempestade alimentada por esse ar uma velocidade de cerca de
\textbf{\qty[text-series-to-math=true]{63}{m.s^{-1}}}, um limite superior.
\end{solution}
